\subsection{Vision--Language Models: What a Recovery Stage Changes}
\label{app:vlm-recovery}

Every number in \cref{sec:exp-vlm} is selection only: a mask is chosen, installed, and scored, and
nothing is trained. This subsection reports what happens when one identical recovery stage is added
to every arm, because it is the comparison a reader of the pruning literature expects and because
what a recovery budget does to a pruned model is worth reporting in its own right.

\textbf{The stage.} LoRA \citep{hu2022lora} of rank $128$ and $\alpha{=}256$ on the language tower, the vision encoder
frozen and the multimodal projector trained full-rank at $2\times10^{-5}$, which is the arrangement
LLaVA's own recovery script uses. The learning rate is LLaVA's $2\times10^{-4}$ scaled linearly to
this run's global batch of $16$, giving $2.5\times10^{-5}$. The corpus is $3000$ image--instruction
samples, half long-form conversation from LLaVA's own instruction data
\citep{liu2023llava} (COCO \emph{train} and Visual Genome images) and half short-answer VQA from
TextVQA \emph{train} \citep{singh2019textvqa} (OpenImages), carrying the same
instruction suffixes the benchmarks are scored under. Both halves are training splits, and no
evaluation image of the seven suites appears in either. The mask stays installed throughout, and every arm of
a cell gets this recipe unchanged --- the only thing that differs between two rows is which units
were removed. Each run re-scores its own pruned model before training and prints the comparison
against the selection-only row it will be compared with, so a recovered number and the number above
it are the same harness, not two.

\textbf{What it does.} On three of the four cells recovery lifts every arm --- by $6.2$ to $21.7$
points of seven-benchmark mean --- and pulls them together: the spread across the four arms falls
from $6.4$ to $2.2$ points on Qwen2.5-VL-7B at $\rho{=}30\%$, from $2.4$ to $0.5$ at $\rho{=}40\%$,
and from $20.1$ to $11.2$ on Qwen3-VL-8B at $\rho{=}30\%$. \method finishes first in two of those
three and second by $1.2$ points in the other; two runs of one configuration on a contended card
differ by $3.2$ points, so that second place is a tie rather than a loss.

The fourth cell behaves differently and is the more informative one. At $\rho{=}40\%$ on
Qwen3-VL-8B the four masks start within $2.0$ points of each other and end $24.2$ apart: LLM-Pruner's
gains $23.9$, ours $15.8$, and Magnitude's gains nothing at all ($-0.3$, its held-out perplexity
still above $90$). Recovery is not a fixed increment added to whatever selection produced. It is a
second optimisation whose reachable endpoint depends on what the mask left behind, and at a budget
where no arm retains a quarter of its dense score above chance before training, that dependence is
what decides the ordering --- not the selection-only scores, which are indistinguishable there.

Two things follow. Where a recovery budget is spent, the ranking it produces is its own: it agrees
with the selection-only ranking on the model that still holds most of its dense score and departs
from it on the one that does not. Selection and recovery are therefore optimising different things at
the highest budget, and mixing them would report neither cleanly --- which is why the body measures
selection alone, the quantity this paper's estimator predicts, and why this table is reported
separately rather than folded into \cref{tab:vlm}. The second reading is the practical one: with
recovery, $\rho{=}40\%$ becomes
usable at all. On Qwen3-VL-8B every recovered arm but Magnitude ends above $46.99$, the best
selection-only score any baseline reaches at $\rho{=}30\%$, and the strongest recovered $40\%$ model
--- LLM-Pruner's, at $59.8$ --- edges past the strongest unrecovered $30\%$ one, ours at $57.7$.

\begin{table}[htbp]
\centering\small
\caption{One identical LoRA recovery stage on every arm, $1000$ items per benchmark.
\emph{sel.\ only} repeats the selection-only row; $+$LoRA is the same model after the stage.}
\label{tab:vlm-recovery}
\adjustbox{max width=\textwidth}{\begin{tabular}{@{}lcccccccc@{}}
\toprule
 & \multicolumn{2}{c}{Qwen2.5-VL-7B, $\rho{=}30\%$} & \multicolumn{2}{c}{Qwen2.5-VL-7B, $\rho{=}40\%$} & \multicolumn{2}{c}{Qwen3-VL-8B, $\rho{=}30\%$} & \multicolumn{2}{c}{Qwen3-VL-8B, $\rho{=}40\%$} \\
\cmidrule(lr){2-3}\cmidrule(lr){4-5}\cmidrule(lr){6-7}\cmidrule(lr){8-9}
Selector & sel.\ only & $+$LoRA & sel.\ only & $+$LoRA & sel.\ only & $+$LoRA & sel.\ only & $+$LoRA \\
\midrule
\textbf{\method} & 63.6 & 69.7 & 44.0 & \textbf{61.9} & 57.7 & \textbf{69.5} & 36.3 & 52.1 \\
LLM-Pruner & 57.2 & \textbf{70.9} & 43.2 & 61.4 & 47.0 & 68.6 & 35.9 & \textbf{59.8} \\
FLAP & 58.5 & 68.7 & 41.6 & 61.5 & 46.5 & 63.9 & 37.9 & 54.5 \\
Magnitude & 58.7 & 69.2 & 41.7 & 61.7 & 37.7 & 58.3 & 35.9 & 35.6 \\
\bottomrule
\end{tabular}
}
\end{table}
